\input{sol_pre.tex}
\input{q_2025-12-29.tex}
\begin{document}
\input{sol_head.tex}

\Savol{1}
$156=2^2\cdot3\cdot13$: \ $\sigma=(1+2+4)(1+3)(1+13)=7\cdot4\cdot14=392$.
\Javob{B}

\Savol{2}
$n^4+n^2+1=(n^2+n+1)(n^2-n+1)$, \ $n^7-n=n(n^3-1)(n^3+1)=(n-1)n(n+1)(n^2+n+1)(n^2-n+1)$, demak har bir qöşiluvçi $\dfrac1{(n-1)n(n+1)}$:
\[\frac1{1\cdot2\cdot3}+\frac1{2\cdot3\cdot4}+\frac1{3\cdot4\cdot5}=\frac{20+5+2}{120}=\frac9{40}.\]
\Javob{B}

\Savol{3}
Birgalikdagi unumdorlik $\frac14+\frac1{12}=\frac13$ iş/kun: $1$ kunda $\frac13$ iş. Qolgan $\frac23$ işni şogird $\frac23:\frac1{12}=8$ kunda bajaradi. Jami $1+8=9$ kun.
\Javob{A}

\Savol{4}
$1\,200\,000\cdot0{,}6\cdot0{,}9=648\,000$.
\Javob{B}

\Savol{5}
$1^3+\dots+59^3=\left(\dfrac{59\cdot60}2\right)^2=1770^2=2^2\cdot3^2\cdot5^2\cdot59^2$. $60=2^2\cdot3\cdot5$ ga bölinadi ($40$, $72$, $48$ uçun $2^3$ yetmaydi).
\Javob{B}

\Savol{6}
$4^9+6^{10}+9^{10}=\left(2^9\right)^2+2\cdot2^9\cdot3^{10}+\left(3^{10}\right)^2=\left(2^9+3^{10}\right)^2$.
\Javob{C}

\Savol{7}
$x=\frac{9-4\sqrt5}2\approx0{,}03$. \ $\sqrt{(x+9)^2}+\sqrt{(x-4\sqrt5)^2}=(x+9)+(4\sqrt5-x)=9+4\sqrt5$.
\Javob{D}

\Savol{8}
$a_{21}+a_{16}=a_{20}+a_{17}=35$, $a_{21}a_{16}=150$ $\Rightarrow$ $\{a_{21};a_{16}\}=\{30;5\}$, $a_{21}-a_{16}=5d$.\\
$a_{16}=5$, $d=5$: $a_2=a_{16}-14d=-65$; \quad $a_{16}=30$, $d=-5$: $a_2=30+70=100$.
\Javob{B}

\Savol{9}
$a_1+a_2+a_3=65-(1+8+35)=21$ $\Rightarrow$ $a_2=7$, $a_1=7-d$, $a_3=7+d$; \ $b_1=8-d$, $b_2=15$, $b_3=42+d$.\\
$15^2=(8-d)(42+d)$ $\Rightarrow$ $d^2+34d-111=0$ $\Rightarrow$ $d=3$ ($d=-37$ da $b_2<b_1$). $b_n$: $5;15;45$, $q=3$.
\[\frac{5(3^n-1)}2=200\Rightarrow3^n=81\Rightarrow n=4.\]
\Javob{A}

\Savol{10}
$\dfrac{a^{-6}\left(1+a^2+a^4\right)}{a^2\left(a^4+a^2+1\right)}=a^{-8}$.
\Javob{A}

\Savol{11}
$ab(a^2+b^2)+(a^2+b^2)+5ab+2(ab)^2=y(x^2-2y)+x^2-2y+5y+2y^2=yx^2+x^2+3y$.
\Javob{D}

\Savol{12}
$\tg2a+\sin2a=\sin2a\cdot\dfrac{1+\cos2a}{\cos2a}=\dfrac{\sin2a\cdot2\cos^2a}{\cos2a}$ $\Rightarrow$ ifoda $=\dfrac{\cos2a}{\sin2a}=\ctg2a$.
\Javob{B}

\Savol{13}
$0\le2\sin x\le1$ $\Rightarrow$ $0\le\sin x\le\frac12$ $\Rightarrow$ $x\in\left[2\pi k;\frac\pi6+2\pi k\right]\cup\left[\frac{5\pi}6+2\pi k;\pi+2\pi k\right]$.
\Javob{C}

\Savol{14}
$17\mp12\sqrt2=(3\mp2\sqrt2)^2=(\sqrt2\mp1)^4$ $\Rightarrow$ $(\sqrt2-1)^x+(\sqrt2+1)^x=6$. \ $t=(\sqrt2+1)^x$: $t+\frac1t=6$ $\Rightarrow$ $t=3\pm2\sqrt2=(\sqrt2+1)^{\pm2}$ $\Rightarrow$ $x=\pm2$. Köpaytma $-4$.
\Javob{D}

\Savol{15}
$x>0$. \ $\ln\dfrac{4(2x+3)}5=\ln\bigl(x(x+3)\bigr)$ $\Rightarrow$ $5x^2+7x-12=0$ $\Rightarrow$ $x=1$ ($x=-\frac{12}5$ mos emas); maxraj $\ln4\ne0$. $1$ ta ildiz.
\Javob{A}

\Savol{16}
Ikkinçi tenglamani $3$ ga köpaytirib birinçisini ayiramiz: $7y=35$, $y=5$, $x=7$. \ $x+y=12$.
\Javob{B}

\Savol{17}
$x^4-3(x-1)^2=\left(x^2-\sqrt3(x-1)\right)\left(x^2+\sqrt3(x-1)\right)$.\\
$x^2-\sqrt3x+\sqrt3=0$: $D=3-4\sqrt3<0$; \quad $x^2+\sqrt3x-\sqrt3=0$: $D=3+4\sqrt3>0$ — $2$ ta ildiz.
\Javob{B}

\Savol{18}
$\dfrac{2x^2+13}{(x-2)(x-3)}+\dfrac{9(x-2)}{(x-2)(x-3)}\le0$ $\Rightarrow$ $\dfrac{(2x-1)(x+5)}{(x-2)(x-3)}\le0$ $\Rightarrow$ $x\in\left[-5;\frac12\right]\cup(2;3)$.\\
Butun yeçimlar: $-5,-4,-3,-2,-1,0$ — $6$ ta.
\Javob{C}

\Savol{19}
Aniqlaniş sohasi $|x|\ge1$. $x<2$ da öng tomon manfiy — tengsizlik bajariladi. $x\ge2$ da: $x^2-1\ge x^2-4x+4$ $\Rightarrow$ $x\ge\frac54$ — bajariladi. Yeçim: $|x|\ge1$. Yeçim bölmaydigan butun son faqat $0$.
\Javob{A}

\Savol{20}
$\left(x+\frac1x\right)^3=x^3+\frac1{x^3}+3\left(x+\frac1x\right)$, \ $f\left(\frac1x\right)=\frac1x+x=f(x)$.
\Javob{C}

\Savol{21}
Bölaklab: $\displaystyle\int_0^\pi(x-1)\cos x\,dx=(x-1)\sin x\Big|_0^\pi-\int_0^\pi\sin x\,dx=0-2=-2$.
\Javob{C}

\Savol{22}
$\lg\sqrt[n]{x+\sqrt{1+x^2}}=\frac1n\lg\left(x+\sqrt{1+x^2}\right)$ — toq funksiya ($(-x+\sqrt{1+x^2})(x+\sqrt{1+x^2})=1$), $\cos(kx^n)$ — juft. Demak $f$ toq, $f'$ juft: $f'(-2023)=2024$, javob $4048$.
\Javob{D}

\Savol{23}
$\dfrac3{\sqrt{x+1}}=2$ $\Rightarrow$ $x+1=\frac94$, $x=1\frac14$.
\Javob{A}

\Savol{24}
Har bir uçburçak töğri burçakli, $O$ dagi burçagi $30^\circ$, keyingi uçburçakning gipotenuzasi oldingisining $O$ ga yopişgan kateti: har qadamda uzunlik $\cos30^\circ$ ga köpayadi. $6$ ta uçburçak:
\[x=64\cos^630^\circ=64\cdot\frac{27}{64}=27.\]
\Javob{B}

\Savol{25}
Balandlik $h=\sqrt3=\frac{\sqrt3}2a$ $\Rightarrow$ $a=2$, \ $S=\frac{\sqrt3}4\cdot4=\sqrt3$.
\Javob{B}

\Savol{26}
$M$ — medianalar kesişgan nuqta, $S_{MAB}=S_{MBC}=S_{MCA}=32$. $M$ medianani $2:1$ böladi:
$MF=\frac23AL-\frac12AL=\frac16AL=\frac14MA$, \ $ME=\frac23BD-\frac13BD=\frac12MB$, \ $MK=\frac23CG'-\frac14CG'=\frac5{12}CG'=\frac58MC$.
\[S_{FEK}=S_{MFE}+S_{MEK}+S_{MKF}=32\left(\frac14\cdot\frac12+\frac12\cdot\frac58+\frac58\cdot\frac14\right)=32\cdot\frac{19}{32}=19.\]
\Javob{B}

\Savol{27}
$EFGH$ töğri törtburçak bölgani uçun kvadrat burçaklaridagi törtta uçburçak öxşaş. $AH=x$, $AE=y$ ($x^2+y^2=a^2$).\\
$\triangle HBG\sim\triangle EAH$, $k=\frac{HG}{EH}=3$: $HB=3y$, $BG=3x$; \ $\triangle FDE\sim\triangle EAH$, $k=\frac{EF}{EH}=3$: $ED=3x$.
\[AB=x+3y=9,\qquad AD=y+3x=11.\]
$x=3$, $y=2$ $\Rightarrow$ $a^2=13$, \ $S=a\cdot3a=3a^2=39$.
\Javob{B}

\Savol{28}
Muntazam beşburçak burçagi $108^\circ$. $DG$ — simmetriya öqi, $\angle ODG=\frac{108^\circ}2=54^\circ$.\\
$\triangle OEF$: $\angle OEF=180^\circ-108^\circ=72^\circ$, $\angle OFE=45^\circ$ $\Rightarrow$ $\angle EOF=63^\circ$.\\
$\triangle ODG$: $\angle DGO=180^\circ-54^\circ-63^\circ=63^\circ=\angle DOG$ $\Rightarrow$ $DG=DO=8$.\\
$EH$ ham simmetriya öqi (uçdan qarama-qarşi tomon örtasiga), barça bunday kesmalar teng: $EH=DG=8$.
\Javob{A}

\Savol{29}
$\vec a(\vec b+\vec c)=\frac12+\frac12=1$, \ $|\vec b+\vec c|^2=1+1+2\cdot\frac12=3$. \ $\cos\varphi=\dfrac1{\sqrt3}$.
\Javob{B}

\Savol{30}
Tekisliklar orasidagi masofa: $h=\sqrt{37^2-12^2}=35$. \ $BD$ proyeksiyasi $\sqrt{125^2-35^2}=\sqrt{14400}=120$.
\Javob{A}

\Savol{31}
1) $x^2-2x-2024=0$: $x=46\in N$ — böş emas. \ 2) $x+19=(x-1)!$: $x=5$ — böş emas.\\
3) $x^2=3\sqrt2-2\sqrt3>0$, $x=\pm\sqrt{3\sqrt2-2\sqrt3}$ — irratsional, böş emas. \ 4) $x=-1$ ildiz — böş emas.\\
5) $|x|=\pi-10<0$ — böş. \quad Faqat $1$ ta.
\Javob{A}

\Savol{32}
Har biri $6$ ta yöldan birini tanlaydi va $5$ km radiusli aylanadagi uçlardan biriga tuşadi. Yöllar orasidagi burçak $60^\circ k$, masofa $2\cdot5\sin(30^\circ k)$: $k=1$ da $5$, $k=2$ da $5\sqrt3$, $k=3$ da $10$. Ikkinçisining yöli birinçisinikiga nisbatan $6$ holatdan $3$ tasida ($k=2,3,4$) masofa $\ge6$:
$P=\frac36=\frac12$.
\Javob{A}

\newpage
\Savol{33}
\fig{\begin{tikzpicture}[scale=0.28,font=\small]
\coordinate (A) at (0,8); \coordinate (H) at (0,0); \coordinate (P) at (6,0); \coordinate (Q) at (15,0);
\fill[okgreen!15] (-15,0)--(A)--(15,0)--(6,0)--(A)--(-6,0)--cycle;
\draw[okgreen] (-15,0)--(A)--(-6,0); \draw[okgreen] (-6,0)--(-15,0);
\draw (A)--(P)--(Q)--cycle; \draw[ax] (A)--(H)--(P);
\draw (0,0.8)--(0.8,0.8)--(0.8,0);
\node[above] at (A) {$A$}; \node[below] at (H) {$H$}; \node[below] at (P) {$B$}; \node[below] at (Q) {$C$};
\node[aux,left] at (0,4) {$8$}; \node[right] at (3.2,4.4) {$10$}; \node[above right] at (8,4) {$17$}; \node[below] at (10.5,0) {$9$};
\end{tikzpicture}}
$p=18$, $S=\sqrt{18\cdot9\cdot8\cdot1}=36$. Eng katta balandlik eng kiçik tomonga ($BC=9$): $h=\frac{72}9=8$.\\
$HB=\sqrt{10^2-8^2}=6$, $HC=\sqrt{17^2-8^2}=15$, $HC-HB=9=BC$ — asos $H$ dan taşqarida. Aylanişda $BC$ halqa çizadi ($r=6$, $R=15$), jism — katta konusdan kiçik konus olib taşlangan.

\textbf{33.} $S_{halqa}=\pi(15^2-6^2)=189\pi\approx567$ — \textbf{D}.

\textbf{34.} Yon sirtlar: $\pi\cdot15\cdot17+\pi\cdot6\cdot10=315\pi\approx945$ — \textbf{B}.

\textbf{35.} $V=\frac13\pi\cdot8\left(15^2-6^2\right)=504\pi\approx1512$ — \textbf{F}.
\Javob{33 — D, \ 34 — B, \ 35 — F}

\newpage
\Savol{36}
$|x|\le1$, \ $x=\cos t$, $t\in[0;\pi]$: \ $\sqrt{1-x^2}=\sin t$, \ $4x^3-3x=\cos3t$:
\[\sin t=\cos3t\ \Leftrightarrow\ \cos\left(\frac\pi2-t\right)=\cos3t\ \Rightarrow\ t=\frac\pi8+\frac{\pi k}2\ \text{ yoki }\ t=-\frac\pi4+\pi k.\]
$[0;\pi]$ da: $t=\frac\pi8,\ \frac{5\pi}8,\ \frac{3\pi}4$ — $3$ ta turli $x$.

a) $3$ ta yeçim.

b) $\cos\frac\pi8\cdot\cos\frac{5\pi}8\cdot\cos\frac{3\pi}4=\cos\frac\pi8\cdot\left(-\sin\frac\pi8\right)\cdot\left(-\frac{\sqrt2}2\right)=\frac{\sqrt2}2\cdot\frac12\sin\frac\pi4=\frac14$.
\Javob{a) 3; \quad b) $\frac14$}

\newpage
\Savol{37}
$c=\cos x$: \ $5(1-c^2)-c^2+8c=|c|-1$ $\Rightarrow$ $6c^2-8c+|c|-6=0$.\\
$c\ge0$: $6c^2-7c-6=0$, $c=\frac{7\pm\sqrt{193}}{12}$ — biri $>1$, biri $<0$, mos emas.\\
$c<0$: $6c^2-9c-6=0$ $\Rightarrow$ $2c^2-3c-2=0$ $\Rightarrow$ $c=-\frac12$ ($c=2$ mos emas).\\
$\cos x=-\frac12$ $\Rightarrow$ $x=\pm\frac{2\pi}3+2\pi k$.

a) $\left[-\frac{5\pi}4;\frac{5\pi}4\right]$ da: $\pm\frac{2\pi}3$ ($\pm\frac{4\pi}3$ çiqib ketadi) — $2$ ta.

b) $[0;6\pi]$ da: $\frac{2\pi}3,\frac{4\pi}3,\frac{8\pi}3,\frac{10\pi}3,\frac{14\pi}3,\frac{16\pi}3$; yiğindi $\frac{54\pi}3=18\pi$.
\Javob{a) 2; \quad b) $18\pi$}

\newpage
\Savol{38}
Juft funksiya: $b=0$, $f(x)=ax^2+c$. Töğri çiziq $Ox$ ni $(1;0)$ da kesadi, demak kesişiş nuqtasi $(1;0)$ va u yagona: $ax^2-4x+c+4=0$ tenglama $x=1$ ikki karrali ildizga ega:
\[x_0=\frac4{2a}=1\Rightarrow a=2;\qquad f(1)=a+c=0\Rightarrow c=-2.\]
$f(x)=2x^2-2$ (tekşiriş: $2x^2-4x+2=2(x-1)^2$).

a) $\dfrac{a+b}c=\dfrac{2}{-2}=-1$.

b) $f(-4)=32-2=30$.
\Javob{a) $-1$; \quad b) 30}

\newpage
\Savol{39}
a) Yarim aylana markazi $Ox$ da, $y_2$ unda: $2\sqrt2x-6\sqrt2=0$ $\Rightarrow$ markaz $(3;0)$. Aylana $(0;0)$ dan ötadi: $R=3$, $x_1=6$.

b) Urinma radiusga perpendikulyar, $y_3\perp y_2$ — demak $y_2$ uriniş nuqtasidan ötadi. $y_2$ ning burçak koeffitsiyenti $2\sqrt2$: $\tg\alpha=2\sqrt2$, $\cos\alpha=\frac13$. Uriniş nuqtasi $T$, markaz $Q(3;0)$, $QT=3$.\\
$\triangle QTX_2$ ($X_2=(x_2;0)$) töğri burçakli ($\angle T=90^\circ$), $\angle TQX_2=\alpha$:
\[QX_2=\frac{QT}{\cos\alpha}=\frac3{1/3}=9\ \Rightarrow\ x_2=3+9=12.\]
\Javob{a) 6; \quad b) 12}

\newpage
\Savol{40}
Grafikdan: eng kiçik qiymat $0$ ($x=0;2\pi;4\pi$), eng kattasi $4$ ($x=\pi;3\pi$). $|a|=\frac{4-0}2=2$, $b=2$, davr $2\pi$ $\Rightarrow$ $|k|=1$. Masalan $f(x)=2-2\cos x$.

a) $\dfrac{|a|+|b|}{|k|}=\dfrac{2+2}1=4$.

b) Böyalgan soha $x=\pi$ dan $x=3\pi$ gaça grafik ostida:
\[S=\int_\pi^{3\pi}(2-2\cos x)\,dx=4\pi-2\sin x\Big|_\pi^{3\pi}=4\pi\approx12.\]
\Javob{a) 4; \quad b) 12}

\newpage
\Savol{41}
Bir uçli, asoslari $AC$ da bölgan uçburçaklar yuzlari asoslariga proporsional.\\
$\frac{S_{BEC}}{S_{ABP}}=\frac{EC}{AP}=\frac49$: $AP=9t$, $EC=4t$. \ $AP+EC=13t=\frac23AC$ $\Rightarrow$ $AC=19{,}5t$, $PE=AC-13t=6{,}5t=\frac13AC$.
\[S_{ABC}=3S_{BPE}=312,\quad S_{ABP}=312\cdot\frac{9t}{19{,}5t}=144,\quad S_{BPC}=312-144=168.\]
$GL$ — örta çiziq, u $BP$ ni ham teng ikkiga böladi ($H$ — $BP$ örtasi).

a) $\triangle BGH\sim\triangle BAP$, $k=\frac12$: \ $S_{HGB}=\frac14\cdot144=36$.

b) $\triangle BHL\sim\triangle BPC$, $k=\frac12$: \ $S_{HLCP}=168-\frac14\cdot168=126$.
\Javob{a) 36; \quad b) 126}

\newpage
\Savol{42}
\fig{\begin{tikzpicture}[scale=0.42,font=\small]
\coordinate (A) at (0,0); \coordinate (D) at (12,0); \coordinate (B) at (0,6); \coordinate (C) at (4,6);
\coordinate (E) at (2,6); \coordinate (F) at (6,0); \coordinate (P) at (3,4.5); \coordinate (C1) at (4,0);
\draw (A)--(B)--(C)--(D)--cycle; \draw (A)--(C); \draw (B)--(D); \draw (E)--(F);
\draw (3,3) circle (3); \draw[ax] (C)--(C1);
\node[below left] at (A) {$A$}; \node[above left] at (B) {$B$}; \node[above] at (C) {$C$}; \node[below right] at (D) {$D$};
\node[above] at (E) {$E$}; \node[below] at (F) {$F$}; \node[right] at (P) {$P$}; \node[below] at (C1) {$C_1$};
\end{tikzpicture}}
Agar diagonallar kesişiş nuqtasi $BC$ va $AD$ örtalarini tutaştiruvçi kesmada yotsa, $AD\parallel BC$ (trapetsiya), va $\triangle PBC\sim\triangle PDA$ da $PE$, $PF$ — mos medianalar: $\frac{AD}{BC}=\frac{PF}{PE}=3$.\\
$BC=b$, $AD=3b$. Içki aylana: $AB+CD=AD+BC=4b$ $\Rightarrow$ $CD=4b-6$.\\
$B$ va $C$ dan $AD$ ga $BB_1$, $CC_1$ balandliklar ($h$), $AB_1=x$, $C_1D=2b-x$:
\[x^2+h^2=36,\quad (x+b)^2+h^2=52,\quad (2b-x)^2+h^2=(4b-6)^2.\]
Birinçisidan: $2bx+b^2=16$; uçinçisidan: $4b^2-4bx+36=16b^2-48b+36$ $\Rightarrow$ $x=12-3b$.\\
$2b(12-3b)+b^2=16$ $\Rightarrow$ $5b^2-24b+16=0$ $\Rightarrow$ $b=4$ ($b=0{,}8$ da $x^2>36$). Demak $x=0$, $AB\perp AD$, $h=6$, $AD=12$, $CD=10$, $C_1D=8$.

a) $\cos\angle ADC=\dfrac{C_1D}{CD}=\dfrac8{10}=\dfrac45$.

b) Aylana ikkala asosga urinadi: $2r=h=6$, $r=3$.
\Javob{a) $\frac45$; \quad b) 3}

\newpage
\Savol{43}
\fig{\begin{tikzpicture}[scale=0.9,font=\small]
\coordinate (A) at (-2.674,0.000); \coordinate (B) at (-1.337,2.316); \coordinate (C) at (1.337,2.316); \coordinate (D) at (2.674,0.000); \coordinate (E) at (1.337,-2.316); \coordinate (F) at (-1.337,-2.316); \coordinate (K) at (-0.344,-1.060);
\coordinate (K1) at ($(D)!1!60:(K)$);
\draw (A)--(B)--(C)--(D)--(E)--(F)--cycle; \draw[okgreen] (B)--(D)--(F)--cycle; \draw (B)--(K)--(D); \draw (F)--(K);
\draw[axs] (D)--(K1)--(F); \draw[ax] (K)--(K1);
\fill (K) circle (1.3pt); \fill (K1) circle (1.3pt);
\node[left] at (A) {$A$}; \node[above left] at (B) {$B$}; \node[above right] at (C) {$C$}; \node[right] at (D) {$D$}; \node[below right] at (E) {$E$}; \node[below left] at (F) {$F$}; \node[above] at (K) {$K$}; \node[below] at (K1) {$K'$};
\end{tikzpicture}}
$BDF$ — muntazam uçburçak, tomoni $t=s\sqrt3$ ($s$ — oltiburçak tomoni).

a) $\triangle DKB$ ni $D$ atrofida $60^\circ$ ga buramiz ($B\to F$, $K\to K'$): $DK'=DK=10$, $\angle KDK'=60^\circ$ $\Rightarrow$ $KK'=10$; $FK'=BK=11$.\\
$\triangle KFK'$: $\cos\varphi=\cos\angle FKK'=\dfrac{25+100-121}{2\cdot5\cdot10}=\dfrac1{25}$, $\sin\varphi=\dfrac{4\sqrt{39}}{25}$. \ $\angle DKF=60^\circ+\varphi$:
\[\cos(60^\circ+\varphi)=\frac1{50}-\frac{\sqrt3}2\cdot\frac{4\sqrt{39}}{25}=\frac{1-12\sqrt{13}}{50},\qquad t^2=DF^2=100+25-100\cos(60^\circ+\varphi)=123+24\sqrt{13}.\]
$s^2=\dfrac{t^2}3=41+8\sqrt{13}$, \ $s=\sqrt{41+8\sqrt{13}}$.

b) $O$ — oltiburçak markazi ($=\triangle BDF$ markazi), $OA=OD=s$, $A$ va $D$ simmetrik. Medianalar formulasi ($KO$ — $\triangle AKD$ medianasi) va muntazam uçburçak uçun Leybnits formulasi:
\[KA^2+KD^2=2\left(KO^2+s^2\right),\qquad KB^2+KD^2+KF^2=3\left(KO^2+s^2\right)=246.\]
$KA^2=\frac23\cdot246-100=64$, \ $AK=8$.
\Javob{a) $\sqrt{41+8\sqrt{13}}$; \quad b) 8}

\newpage
\Savol{44}
\fig{\begin{tikzpicture}[scale=1.5,font=\small]
\coordinate (O1) at (0,0); \coordinate (O2) at (0,2); \coordinate (M1) at (1.155,0); \coordinate (M2) at (0.866,2); \coordinate (Q) at (0,1);
\draw (M1)--(M2); \draw (O1)--(M1); \draw (O2)--(M2); \draw[dashed] (O1)--(O2);
\draw[okgreen] (Q) circle (1);
\draw[ax] (Q)--(M1); \draw[ax] (Q)--(M2);
\fill (Q) circle (1pt);
\node[below] at (O1) {$O_1$}; \node[above] at (O2) {$O_2$}; \node[below right] at (M1) {$M_1$}; \node[above right] at (M2) {$M_2$}; \node[left] at (Q) {$Q$};
\node[below] at (0.58,0) {$r_1$}; \node[above] at (0.43,2) {$r_2$};
\end{tikzpicture}}
Kesim: öq va yon yoq apofemasi orqali. $O_1M_1=r_1=\frac{4}{2\sqrt3}=\frac2{\sqrt3}$, \ $O_2M_2=r_2=\frac3{2\sqrt3}=\frac{\sqrt3}2$ (asoslarga içki çizilgan aylana radiuslari). Şar markazi $Q$ öqda, $QO_1=QO_2=R$, şar $M_1M_2$ ga urinadi. $QM_1$, $QM_2$ — $M_1$, $M_2$ dagi burçak bissektrisalari, ular yiğindisi $180^\circ$ bölgan burçaklarni böladi $\Rightarrow$ $\angle M_1QM_2=90^\circ$. Uriniş nuqtasi $T$ uçun $M_1T=r_1$, $M_2T=r_2$, $QT^2=M_1T\cdot M_2T$:
\[R^2=r_1r_2=\frac2{\sqrt3}\cdot\frac{\sqrt3}2=1\Rightarrow R=1,\quad h=2R=2.\]
a) $R=1$.

b) $S_1=\frac{\sqrt3}4\cdot16=4\sqrt3$, $S_2=\frac{9\sqrt3}4$, $\sqrt{S_1S_2}=3\sqrt3$:
\[V=\frac h3\left(S_1+S_2+\sqrt{S_1S_2}\right)=\frac23\cdot\frac{16\sqrt3+9\sqrt3+12\sqrt3}4=\frac{37\sqrt3}6.\]
\Javob{a) 1; \quad b) $\frac{37\sqrt3}6$}

\newpage
\Savol{45}
\fig{\begin{tikzpicture}[scale=0.55,font=\small]
\coordinate (B) at (0,0); \coordinate (A) at (0,6); \coordinate (C) at (6,0); \coordinate (D) at (6,6); \coordinate (O) at (2,3);
\fill[okgreen!15] (C)--(B)--(0,4.5)--cycle;
\draw (A)--(B)--(C)--(D)--cycle; \draw[axs] (C)--(0,4.5); \draw[ax] (2,-0.3)--(2,6.3);
\draw[dashed] (O)--(0,3); \draw[dashed] (O)--(2,0);
\fill (O) circle (2pt);
\node[below left] at (B) {$B$}; \node[above left] at (A) {$A$}; \node[below right] at (C) {$C$}; \node[above right] at (D) {$D$}; \node[above right] at (O) {$O$};
\node[left] at (0,4.5) {$M$}; \node[above] at (1,3) {\scriptsize$2$}; \node[right] at (2,1.5) {\scriptsize$3$};
\end{tikzpicture}}
$O$ dan $AB$ gaça $2$, $BC$ gaça $3$ m.

a) Töğri çiziq kvadrat burçagidan uçburçak kesadi. $O$ dan ötuvçi çiziqni burganda uçburçak yuzi kateti kvadrat tomoniga ($6$ m) yetgunça ösadi, ya'ni eng katta uçburçak çiziq kvadrat uçidan ötganda hosil böladi:\\
$C$ va $O$ orqali: $AB$ ni $M$ da kesadi, $\frac{BM}{3}=\frac{6}{4}$ $\Rightarrow$ $BM=4{,}5$, \ $S=\frac12\cdot6\cdot4{,}5=13{,}5$;\\
$D$ va $O$ orqali ham $13{,}5$; $A$ yoki $B$ orqali: $\frac12\cdot6\cdot4=12$. \ Eng kattasi $13{,}5$ m$^2$.

b) $|S_1-S_2|=36-2S_{\min}$, eng kiçik bölakni topamiz.\\
1) Çiziq $AD$ va $BC$ ni kessa: $AB$ tomondagi bölak — trapetsiya, balandligi $6$, örta çiziği $O$ dan ötadi va $2$ ga teng: $S=12$ (çiziqqa boğliq emas).\\
2) Çiziq $AB$ va $CD$ ni kessa: $BC$ tomondagi trapetsiyaning örta çiziği $3\pm$ (köpi bilan $0{,}75$), yuzi $13{,}5$ dan $22{,}5$ gaça.\\
3) Burçak uçburçaklari: $O$ gipotenuza örtasida bölganda eng kiçik, $B$ va $A$ burçaklarida $\frac12\cdot4\cdot6=12$, $C$ va $D$ da undan katta.\\
$S_{\min}=12$ m$^2$ $\Rightarrow$ $|S_1-S_2|_{\max}=36-24=12$ m$^2=120\,000$ sm$^2$.
\Javob{a) 13,5; \quad b) 120\,000}

\ShowJavoblar
\end{document}
